\title{Large Language Models Can Automatically Engineer Features for Few-Shot Tabular Learning}

\begin{abstract}
Large Language Models (LLMs) possess impressive capabilities for solving complex, unfamiliar reasoning tasks, which positions them as powerful tools for tabular data analysis—a domain integral to many practical settings. We introduce a new in-context learning approach named \textsf{FeatLLM}, which leverages LLMs to act as feature constructors, thus crafting a transformed input dataset that enhances the effectiveness of tabular prediction tasks. These engineered features are subsequently employed by a lightweight downstream machine learning model, such as linear regression, which enables highly effective few-shot learning. Notably, once the optimal features have been identified, \textsf{FeatLLM} utilizes only this basic predictive model during inference, omitting further interactions with the LLM. Unlike previous LLM-centric strategies, \textsf{FeatLLM} removes the necessity to interact with the LLM for each individual test instance at inference. Furthermore, the approach relies solely on API-level connectivity with LLMs and circumvents constraints posed by prompt length. Experiments spanning a variety of tabular datasets and domains reveal that \textsf{FeatLLM} uncovers robust predictive rules, achieving on average a 10\% improvement over alternatives such as TabLLM and STUNT.
\end{abstract}

\section{Introduction}
The last several years have witnessed significant advances by Large Language Models (LLMs) in generating accurate outputs for previously unseen tasks, all while avoiding the need for dedicated fine-tuning for each task~\cite{creswell2022selection,imani2023mathprompter,taylor2022galactica}. Thanks to pre-training on vast text datasets, LLMs encapsulate broad prior knowledge, which can often replace human expert input in numerous fields~\cite{Genomes2021,singhal2023large,wilcox2003role}, and have become foundational for automating a diverse array of applications, such as dialogue systems~\cite{finch2023leveraging} and automatic program repair~\cite{fan2023automated}. By designing prompts that include both the specifics of the task and several illustrative examples, these models have the ability to grasp context and pay selective attention to pertinent attributes and cases~\cite{brown2020language,zhao2023survey}. Such proficiency showcases LLMs’ ability to interpret and process data, indicating their suitability as advanced feature engineers.

The extension of LLMs’ reasoning and general knowledge to tabular learning has garnered substantial research attention. To illustrate, prior knowledge—for example, the fact that advancing age correlates with increased risk of specific diseases—can enhance models tasked with predicting health outcomes. Contemporary studies often represent feature and task information in natural language and serialize data prior to submission to an LLM~\cite{dinh2022lift,wang2023anypredict}. Subsequent processing may employ in-context learning~\cite{nam2023semi} or apply lightweight fine-tuning methods~\cite{dinh2022lift,hegselmann2023tabllm}. Integrating LLM-derived prior knowledge has demonstrated effectiveness, yielding superior performance compared to traditional tabular baselines, particularly in settings with limited labeled data~\cite{hegselmann2023tabllm}.

Current tabular learning approaches that leverage LLMs encounter several notable drawbacks. In scenarios requiring end-to-end prediction, performing inference necessitates running the LLM at least once for every data point, which incurs considerable computational overhead. Achieving high predictive accuracy typically demands that the LLM undergo fine-tuning; this dependency significantly hinders deployment, especially when working with LLMs lacking accessible source code or tuning endpoints. This challenge is exacerbated by the fact that many state-of-the-art LLMs now restrict access solely to inference APIs~\cite{anil2023palm,openai2023gpt4}. Moreover, most current techniques are poorly suited to tasks involving extensive prompts~\cite{liu2023lost}: as the feature dimensionality of the dataset increases and the corresponding textual representation grows in length, these methods often either become computationally impractical or suffer notable drops in performance. Collectively, these issues create substantial obstacles to the practical adoption of LLM-powered tabular learning in real applications.

To address these challenges, our method repositions the role of LLMs from directly producing predictions to serving as tools for feature construction. By capitalizing on their inherent domain knowledge, we enable LLMs to participate in efficient feature engineering. We put forth a new in-context learning paradigm, termed \textsf{FeatLLM}, designed to probe and extract the “criteria” that inform LLM decisions. In practice, for a given classification problem—such as disease prediction—the LLM is tasked with deducing and articulating rule-like conditions based on input features that determine positive cases. The extracted criteria then provide the basis for generating new, informative features that supplant existing ones in subsequent modeling steps. This repositioning makes downstream inference substantially simpler: sophisticated machine learning models are no longer necessary after relevant features are engineered, which reduces inference time. Leveraging LLMs’ ability for in-context learning, we avoid model retraining entirely by relying on prompt-based demonstrations to extract such features. To further contend with the limitation imposed by lengthy inputs, we incorporate ensemble-inspired methods~\cite{breiman2001random} such as feature bagging, thus effectively controlling prompt size.

\textsf{FeatLLM} decomposes the process of constructing input prompts for feature engineering into two sequential sub-tasks: (1) comprehending the underlying problem to infer connections between features and target variables; and (2) utilizing the previously acquired insights, along with a handful of illustrative examples, to articulate critical decision rules that facilitate class differentiation. By guiding the LLMs through this organized pattern of reasoning, the framework encourages exclusion of irrelevant features during rule formulation. The formulated rules are subsequently applied to the data (executed as program code), resulting in a set of binary features, each indicating whether an instance satisfies a particular rule. Subsequently, a simple model is trained on these engineered features to predict class probabilities. To further bolster robustness, an ensemble method is adopted in which feature generation is repeated with feature bagging. The use of adjustable numbers of features and sample sizes during bagging effectively addresses concerns about excessive prompt lengths. Since \textsf{FeatLLM} operates without requiring training and incurs minimal inference costs, it renders LLM-driven feature engineering practical for a broad range of real-world scenarios.

\textsf{FeatLLM} is empirically assessed across 13 different tabular datasets under low-shot conditions, demonstrating consistently robust and superior performance. The results highlight the LLMs’ capacity to integrate both pre-existing domain knowledge and empirical evidence from data, resulting in the synthesis of informative features. Compared to current few-shot learning methods, our framework consistently achieves better results in multiple scenarios. The implementation is publicly available at the anonymized GitHub repository~\texttt{https://github.com/Sungwon-Han/FeatLLM}.

\section{Related Work}

\subsection{Few-Shot Learning with Tabular Data}

Few-shot learning has seen notable progress, enabling the derivation of transferable representations from datasets with only a minimal amount of labeled examples~\cite{chen2018closer}. While its initial successes were concentrated in computer vision, research has subsequently demonstrated the adaptability of few-shot learning to other formats, such as text, audio, and, more recently, tabular datasets~\cite{majumder2022few,nam2023stunt,schick2021exploiting}. The primary difficulty with learning from small labeled datasets is the heightened susceptibility to overfitting and learning spurious associations~\cite{bartlett2021deep}. To address this, new lines of research have augmented labeled data with supplementary sources—either by injecting domain-specific priors or integrating auxiliary information—which are intended to induce appropriate inductive biases. One prominent direction leverages the abundance of unlabeled data, employing creative data augmentation techniques to improve semi-supervised learning outcomes~\cite{ucar2021subtab,yoon2020vime}; another strives for robust generalization through task-independent unsupervised pre-training that provides strong model initializations~\cite{somepalli2021saint}. These unsupervised schemes often include reconstructive approaches, such as masking or intentionally corrupting data then training the model to recover it~\cite{arik2021tabnet,majmundar2022met}, or using data augmentation in tandem with contrastive learning strategies~\cite{bahri2022scarf,chen2020simple}. Nevertheless, the inherent heterogeneity of tabular datasets poses significant difficulties for designing augmentation protocols that are universally effective, and empirical results indicate these methods typically yield only modest improvements for few-shot learning on tabular data~\cite{nam2023stunt}.

In response, current trends have explored pre-training models on a diverse collection of tabular tasks—spanning beyond the immediate application—and subsequently adapting or transferring the learned knowledge to specific target tasks~\cite{zhang2023generative,levin2022transfer,zhu2023xtab}. For example, TransTab~\cite{wang2022transtab} applies transformer-based architectures to infer semantic linkages across columns by interpreting column names as context, rather than just the data entries themselves. This allows the system to repurpose knowledge across tables, resulting in improved performance for both zero-shot and few-shot predictions on unfamiliar tasks. Alternatively, TabPFN~\cite{hollmann2022tabpfn} introduces a different strategy, generating synthetic tabular data from an assortment of distributions to pre-train a Bayesian neural network. After this stage, the model can produce inferences for a target task simply by ingesting both training and test samples simultaneously, bypassing the need for further training. The broad statistical understanding gained in this pre-training phase consistently surpasses classical tree-based algorithms and shows particular promise in the low-data regime.

\subsection{Language-Interfaced Tabular Learning}

Leveraging large language models (LLMs)~\cite{anil2023palm,openai2023gpt4} provides another compelling direction for incorporating prior knowledge into tabular learning. Due to their training on expansive and diverse text datasets, LLMs exhibit remarkable capabilities for generalization, often successfully applying knowledge to novel tasks across different domains~\cite{brown2020language}. Recent research has investigated ways to capitalize on LLMs’ capacity for broad knowledge, specifically targeting tabular data tasks. The standard approach involves converting, or serializing, the structured tabular information into textual representations, which can then be introduced as input prompts to LLMs~\cite{dinh2022lift,hegselmann2023tabllm,wang2023anypredict}. Through careful construction of these prompts—including explicit inclusion of tabular contents, explanatory feature labels, and detailed task descriptions—the models are able to reason about relationships among tasks and table attributes to generate predictions. Innovations in this space extend to the application of parameter-efficient tuning methodologies such as LoRA~\cite{hu2021lora} and IA3~\cite{liu2022few}, which allow specialized training for these models, as well as the usage of in-context learning—where a handful of training examples are embedded in the prompt to elicit desired behaviors from the LLMs, all without explicit parameter updates~\cite{dinh2022lift,hegselmann2023tabllm,nam2023semi,slack2023tablet}.

Although the strategies discussed previously have shown efficacy in enhancing few-shot tabular learning, their reliance on channeling every inference through the LLM introduces obstacles for real-world industrial deployment. The present work seeks to shift away from the standard black-box paradigm, where each individual prediction is routed through an LLM in a holistic manner. Rather, our approach prompts the LLM to derive the underlying conditions or rules that define each target class directly. \textsf{FeatLLM} subsequently converts these inferred rules into executable code that engineers additional feature columns. As a result, predictions can be made without recurring to the LLM, harnessing in-context learning to derive rules by leveraging both prior knowledge and few-shot exemplars.

\section{Methods}
\textsf{FeatLLM} is tailored to extract the class-specific “rules”—criteria that describe the assignment to each label—from a set of few-shot examples, as illustrated in Figure~\ref{fig:main_model}. The rules inferred from an LLM are then programmatically parsed to generate new binary features for each sample, where each feature indicates whether a specific rule is met. These binary indicators are aggregated through a dot product with a set of trainable non-negative weights, resulting in the estimated probability for each possible class. Training this framework allows the model to calibrate the significance of each binary rule feature and thus refine class likelihood predictions (see Section~\ref{Sec:training}). This entire workflow is executed multiple times with bagging, and outcomes are consolidated via ensembling. The formulation of the problem and detailed discussion of each procedural step are provided in the subsequent subsections.

\vspace*{-0.12in}
\paragraph{Problem formulation.} We examine a tabular dataset comprising $N$ labeled instances, denoted as $\mathcal{D} = \{(\mathbf{x}^i, \mathbf{y}^i)\}_{i=1}^{N}$. Each sample $\mathbf{x}^i$ corresponds to a $d$-dimensional input vector, representing $d$ input features. Feature names, written in natural language, are collected as $F = \{f_j\}_{j=1}^{d}$, and their respective definitions are provided elsewhere. Assuming a classification scenario, each label $\mathbf{y}^i$ is an element from a set of predefined categories. In $k$-shot learning setups, the model is trained using only $k$ labeled examples, selected at random where $k < N$.

\subsection{Prompt Design for Extracting Rules}
\label{Sec:prompt_design}
To encourage the LLM to extract rules by simulating a logical reasoning process, we structure the prompt to reflect the steps an expert might take when tackling a tabular learning problem. Our prompt is structured into three primary sections (refer to Fig.~\ref{fig:prompt_for_rule} and Appendix~\ref{sec:example_rule}):

\vspace*{-0.12in}
\paragraph{Basic information description.} This initial segment presents key data necessary for addressing the problem (highlighted in orange in Fig.~\ref{fig:prompt_for_rule}). Here, the problem statement—with associated labels—is given, as well as descriptions of each feature and illustrative examples. The problem statement takes the form of a question (e.g., ``Does this patient's myocardial infarction complications data indicate chronic heart failure? Yes or no?'' See Appendix~\ref{sec:task_desc} for further task descriptions.). Each feature is specified with its type of value, and definitions may be supplied when available. In cases involving categorical variables, possible category values can also be illustrated. For demonstration purposes, a few training instances are transformed into text, explicitly paired with their actual labels. The process of text serialization is adapted from prior work~\cite{dinh2022lift,hegselmann2023tabllm} and formatted as:

\begin{align}
&\text{Serialize}(\mathbf{x}^i,\mathbf{y}^i, F) = \nonumber \\
&``\ f_1\ \text{is}\ \mathbf{x}^i_1.\ ... \ f_d\  \text{is}\  \mathbf{x}^i_d.\ \ \text{Answer:}\  \mathbf{y}^i\ ”
\end{align}

\vspace*{-0.12in}
\paragraph{Reasoning instruction.} Our objective is to improve the reasoning capability of the LLM by explicitly directing its reasoning process through reasoning instructions (refer to the blue highlights in Fig.~\ref{fig:prompt_for_rule}). The reasoning instruction starts with an opening statement, drawing inspiration from the chain-of-thought paradigm~\cite{wei2022chain}. Subsequently, the LLM’s reasoning is structured into a two-phase procedure. During the initial phase, the LLM is prompted to discern possible causal links or tendencies between features and the overarching task description. This inference is performed independently of provided demonstrations, leveraging only established knowledge or intuitive reasoning, which helps maximize the activation of the model’s prior understanding of the issue. In the following phase, the LLM incorporates both the insights gained in the first phase and the provided example demonstrations to generate rules pertinent to each target class. Segmenting the reasoning in this manner helps mitigate the model’s susceptibility to spurious associations among irrelevant variables and enables it to prioritize more pertinent features. To regulate the comprehensiveness of rule extraction and to strike a balance between underfitting and unwieldy prompt length, we prescribe a fixed extraction count per class, specifically 10 rules\footnote{Further discussion on rule quantity is available in Section~\ref{sec:analysis}.}. Moreover, in the rule construction process, the LLM is prompted to consider the feature type provided in the initial information summary to prevent mismatches between rules and data types. 

\vspace*{-0.12in}
\paragraph{Response instruction.} To optimize the interpretability and downstream application of the inferred rules, the LLM is given explicit directions on how to present its output (see yellow highlights in Fig.~\ref{fig:prompt_for_rule}). The prompt defines both the required response schema for each answer category (i.e., the prescribed response format) and the template that each rule should conform to (i.e., the specified rule structure). This guidance assists the LLM in appropriately tailoring rule formats according to feature types. In the absence of additional directions, the LLM may opt to compose rules that integrate conditions with logical connectors, such as AND and OR. Illustrative samples of expected responses are detailed in Appendix~\ref{sec:outcome-rule}.

\subsection{Inferring Class Likelihood via Rules}
\label{Sec:training}

\paragraph{Parsing rules for feature generation.} The rules produced in the earlier phase are used to construct novel binary features. For each class, we derive features that represent whether a sample fulfills the class-specific rules (where the outcome is ‘0’ or ‘1’). These binary indicators can serve as informative variables for assessing class membership likelihood. As the rules are articulated in natural language by the LLM, it becomes necessary to establish a robust parsing method to automate the process of feature extraction. Although the response and rule formats are standardized to facilitate parsing, LLM-generated output often exhibits inconsistent or unexpected syntactic modifications~\cite{kovavc2023large}. For example, instead of using the symbols ``$>$” or ``$<$”, the model might produce equivalent verbal expressions, such as ``greater than" or ``less than", which increases parsing difficulty.

To effectively handle the noisy and variable nature of text outputs, rather than relying on intricate, hardcoded parsers, we instead employ the LLM’s inherent semantic parsing abilities. The LLM excels at inferring intent even when syntax is irregular, and can generate simple code following explicit instructions—an ability underpinned by its training on source code. To harness these strengths, our prompt to the LLM specifies the target function’s name, a description of both inputs and outputs, and the rules to be implemented. We refer the reader to Figures~\ref{fig:prompt_for_parsing} and~\ref{sec:example_parsing}-\ref{sec:example_parse_outcome} in the Appendix for exemplar prompts and resultant outputs. The generated code is executed using Python’s \texttt{exec()} method to perform the required data transformation.

\vspace*{-0.12in}
\paragraph{Inferring class likelihood.} When new binary features are produced from class-specific rules, a straightforward approach for estimating the likelihood that a sample belongs to a particular class is by tallying the number of rules for that class satisfied by the sample—that is, summing the corresponding binary features within each class. However, since individual rules and features may vary in their discriminative power, it is important to quantify their relative contributions by leveraging information from training data. To address this, we propose to model feature importance through a standard machine learning (ML) framework that operates on the set of binary features associated with each class. As an example, we may employ a bias-free linear model. Specifically, let the binary features for class $k$ and sample $\mathbf{x}^i$ be expressed as $\mathbf{z}_k^i \in \{0, 1\}^R$, with $R$ denoting the number of rules per class and $c$ the total number of classes. The prediction probability vector $\mathbf{p}^i$ for class membership can then be obtained as follows:

\begin{align}
\text{logit}_k^i &= \max(\mathbf{w}_k, 0) \cdot \mathbf{z}_k^i, \nonumber \\
\mathbf{p}^i &= \text{Softmax}({[\text{logit}_{1}^i, ..., \text{logit}_{c}^i]}). \label{eq:infer_prob}
\end{align}

In this formulation, $\mathbf{w}_k$ stands for the vector of learnable coefficients within the linear model, each encoding the weight or relevance of its respective feature. By enforcing that the learned weights remain non-negative through projection onto the positive orthant, we prevent class-supporting features induced by the LLM from inadvertently reducing the class score during training. Ultimately, the computed logit vector is normalized to class probabilities using the softmax function. To estimate $\mathbf{w}_k$, we minimize the cross-entropy loss over a limited set of labeled examples (few-shot learning).

\vspace*{-0.12in}
\paragraph{Ensembling with bagging.} To increase the reliability of predictions, we replicate the aforementioned methodology multiple times, thereby generating an ensemble of independently trained inference models. For each trial $t$ ($t=1, ..., T$), let $\mathbf{w}[t]$ denote the resulting weight vector. The final prediction for an input is then determined by averaging the class probability vectors $\mathbf{p}^i$, each computed as described in Eq.~\ref{eq:infer_prob} with a distinct weight set $\mathbf{w}[t]$.

To inject variability into rule generation for ensemble methods, we set the LLM inference temperature to a high value. Additionally, we shuffle the sequence of few-shot demonstrations, motivated by findings that changing their order can influence LLM responses~\cite{lu2022fantastically}\footnote{Detailed analyses of rule diversity are provided in Appendix~\ref{sec:diversity}}. In order to leverage the diversity yielded in the prediction stage for enhanced robustness, we employ a bagging approach whereby, for each ensemble trial, we sample a subset of features or instances—a strategy that becomes increasingly beneficial when the number of features or training examples is substantial. The ensemble strategy we introduce yields two chief benefits. Firstly, it increases the resilience of the overall framework: if the LLM produces incorrect rules in particular trials, accurate rules generated in other ensemble trials can compensate for these, thereby boosting collective robustness. This principle aligns with the self-consistency property of LLMs~\cite{wang2022self}, as rules that are repeatedly discovered across ensemble iterations are more likely to be valid. Secondly, the ensemble approach mitigates LLM prompt length constraints. For large tabular datasets that would otherwise exceed the LLM’s permissible prompt size, bagging enables working with manageable input subsets, sustaining robust model performance. While a single rule set may not adequately represent the influence of all features, aggregating multiple weaker models, each drawn from different combinations of features or instances, helps to counteract limitations of feature/instance coverage in any one ensemble trial.

\section{Experiments}
We assess \textsf{FeatLLM} on a range of tabular data benchmarks and undertake several analyses to better understand our framework’s effectiveness. Further experiments, encompassing different LLM models, qualitative assessments, and more, are detailed in the Appendix due to space limitations.

\subsection{Performance Evaluation}
\paragraph{Datasets.}
Our empirical study utilizes 13 publicly available datasets for tasks involving binary or multi-class classification: (1) Adult~\cite{asuncion2007uci}, where the aim is to classify whether an individual earns more than \$50,000 per year; (2) Bank~\cite{moro2014data}, used to predict if a client will subscribe to a term deposit; (3) Blood~\cite{yeh2009knowledge}, which involves forecasting whether a blood donor will return for future donations; (4) Car~\cite{kadra2021well}, for categorizing car quality; (5) Communities~\cite{misc_communities_and_crime_183}, used for predicting the crime rate in specific localities; (6) Credit-g~\cite{kadra2021well}, to classify whether an individual poses a good or bad credit risk; (7) Diabetes\footnote{\texttt{kaggle.com/datasets/uciml/pima-indians-diabetes-database}}, aiming to detect diabetes incidence; (8) Heart\footnote{\texttt{kaggle.com/datasets/fedesoriano/heart-failure-prediction}}, designed for coronary artery disease prediction; and (9) Myocardial~\cite{misc_myocardial_infarction_complications_579}, to determine whether an individual experiences chronic heart failure.

Additionally, we have incorporated both recently released and synthetic datasets that are relatively understudied and were absent from the pre-training corpus of large language models. Specifically, these include: (10) Cultivars, aimed at forecasting the grain yield for a particular soybean cultivar\footnote{\texttt{archive.ics.uci.edu/dataset/913}}; (11) NHANES, focused on predicting the age group of an individual based on available records\footnote{\texttt{archive.ics.uci.edu/dataset/887}}; (12) Sequence-type, where the task is to classify sequences into one of four categories: arithmetic, geometric, Fibonacci, or Collatz; and (13) Solution-mix, which involves determining if a mixture created from four solutions has a percentage concentration greater than 0.5. The first two datasets were contributed to the UCI Machine Learning Repository after September 2023, ensuring their exclusion from the training data of the LLM API (gpt-3.5-turbo-0613) adopted in our experiments. The last two datasets are artificially constructed, guaranteeing they were not present in the LLM API’s pre-training set, thereby validating their use for robust model evaluation.

Table~\ref{Tab:basic-desc} presents a summary of the diverse sizes and complexities found in these datasets. Each one offers well-defined attributes, with clear names and descriptive information. Notably, datasets like `Communities' and `Myocardial' contain more than 100 features, which can create difficulties due to the limited input length manageable by LLMs.

\vspace*{-0.12in}
\paragraph{Baselines.}
Our experimental evaluation contrasts the proposed approach against nine baselines. The first set comprises three established tabular learning techniques: (1) Logistic regression (LogReg), (2) XGBoost~\cite{chen2016xgboost}, and (3) RandomForest~\cite{ho1995random}. The subsequent pair involves methods that leverage unlabeled datasets: (4) SCARF~\cite{bahri2022scarf} and (5) STUNT~\cite{nam2023stunt}. It is important to acknowledge that access to extensive unlabeled datasets is often impractical outside of laboratory settings. We also include (6) TabPFN~\cite{hollmann2022tabpfn} as a baseline, which relies on pretraining the model over diverse, synthetically generated datasets. Finally, three LLM-based baselines are considered: (7) In-context learning (In-context)~\cite{wei2022emergent}, (8) TABLET~\cite{slack2023tablet}, and (9) TabLLM~\cite{hegselmann2023tabllm}. In-context learning provides examples directly in the prompt for prediction, without model fine-tuning. TABLET supplements the prompt with additional interpretability cues via rule sets and prototype samples, integrating these using an external classifier to support improved inference. TabLLM implements parameter-efficient adaptation, for instance using IA3~\cite{liu2022few}, allowing few-shot training of the LLM on the target dataset.

\vspace*{-0.12in}
\paragraph{Implementation details.}
\textsf{FeatLLM} is primarily implemented with GPT-3.5 as the foundational LLM, although the approach remains flexible with respect to the underlying LLM architecture (refer to Appendix~\ref{sec:backbones} for analysis on alternative backbones). LLM inference is performed using a temperature parameter of 0.5, while the top-p value is retained at its API default setting of 1. For our method, we configure the number of ensemble members to 20, and limit the extracted rules to 10. Further analysis concerning the selection of hyper-parameters is shown in Figure~\ref{fig:hyperparameter2} and discussed in Appendix~\ref{sec:hyper-parameter}.

The linear models are trained using the Adam optimizer with a learning rate set at 0.01 for up to 200 epochs. Epoch selection is carried out through $k$-fold cross-validation. When reproducing baseline results, we employ GPT-3.5 for in-context learning scenarios. For experiments with TabLLM, T0~\cite{sanh2022multitask} is used in alignment with original experiment configurations, as TabLLM limits LLM choices to those with publicly released checkpoints, making GPT-3.5 unavailable for that benchmark. For conventional machine learning algorithms such as LogReg, XGBoost, and RandomForest, hyper-parameters are chosen via Grid-search alongside $k$-fold cross-validation. The value for $k$ is assigned as either 2 or 4, ensuring all classes are present in the training partition. Additional implementation details can be found in Appendix~\ref{sec:baselines}.

\vspace*{-0.12in}
\paragraph{Results.}
Tables~\ref{tab:main1} and \ref{tab:main2} display comparative results of model performance across various datasets. All experimental evaluations were conducted in triplicate, each with different random partitions of the training and test sets; mean and standard deviation for the AUC metric are reported. \textsf{FeatLLM} persistently emerges as either the best-performing or second-best model in relation to the baseline competitors. Impressively, our framework matches the performance of conventional tabular learning methods operating with 32 training examples, even when only 4 examples are provided for adaptation (refer to Fig.~\ref{fig:compares}). Although STUNT and TabPFN attained notable improvements on specific datasets, their overall gains relative to standard machine learning approaches were modest. In addition, LLM-based models like in-context learning, TABLET, and TabLLM encountered inference failures on tabular tasks with high feature dimensionality. In contrast, our method, which integrates bagging and ensembling strategies, maintained robust results even with datasets encompassing numerous features. Our bagging and ensemble methodology also offers potential for integration into existing LLM-based models; however, such incorporation would necessitate approximately doubling inference passes per sample, thus substantially increasing the computation requirements and limiting practicality.

\subsection{Analysis \& Discussion}
\label{sec:analysis}
\paragraph{Ablation study.} 
To assess the contributions of key methodological elements, we perform ablation studies considering: (1) \textsf{-Tuning}, where the linear model’s weight adjustment is omitted and logits are derived by summing binary feature indicators; (2) \textsf{-Ensemble}, where the ensemble mechanism is removed; (3) \textsf{-Description}, where the provision of feature descriptions is excluded; and (4) \textsf{-Reasoning}, where Step-1 of the rule-generation reasoning instruction is withheld.

Table~\ref{tab:ablation} reports the AUC differences resulting from these ablations, averaged across datasets. For every ablated variant, a reduction in model performance is observed. Notably, omitting weight tuning has an increasingly detrimental effect as the number of training examples grows, implying that, with more data, correctly weighting the rules becomes more advantageous. Conversely, omitting feature descriptions or reasoning instructions has the most pronounced negative effect in the low-shot regime, highlighting the value of leveraging the prior knowledge embedded in LLMs when limited training data is available. Ultimately, the inclusion of the proposed ensemble consistently delivers performance benefits across all examined scenarios.

\vspace*{-0.12in}
\paragraph{Training \& inference time comparison.} To assess computational efficiency, we compare both training and inference durations across the examined methods. All experiments are carried out on the Adult dataset using a training subset consisting of 16 samples. A single A100 GPU serves as the standard hardware configuration for all models except TabLLM, which employs four A100 GPUs to facilitate model parallelism. Table~\ref{tab:cost} summarizes the execution time metrics. Of particular note, \textsf{FeatLLM} achieves low inference latency, on par with traditional baselines. By contrast, alternative approaches relying on LLMs—including In-context learning, TABLET, and TabLLM—require considerably longer inference times. As for training, \textsf{FeatLLM} demands a number of API calls (just 30), a quantity similar to other LLM-based techniques, which neither significantly increases costs nor precludes parallel execution.

\vspace*{-0.12in}
\paragraph{Quality of features generated by \textsf{FeatLLM}.}
To evaluate the practical utility of the rule-based features extracted by \textsf{FeatLLM}, we perform a straightforward analysis. For each dataset, we compute the mean AUROC of the newly constructed features against the target labels, subsequently identifying the proportion of features surpassing an AUROC of 0.5. Table~\ref{tab:quality} reports these findings. The data clearly indicate that most generated rules capture signal relevant to the task, contributing valuable information to the predictive target (see the column ``Before tuning"). Additionally, during the training of the linear model, the method inherently excludes features with limited relevance—specifically, those assigned weights below a specified threshold (displayed in the ``After tuning" column). In this case, the filtering threshold is established at 0.1, based on the empirical distribution of learned weights.

\vspace*{-0.12in}
\paragraph{Effect of adding spurious correlations.} To evaluate how different techniques mitigate the influence of spurious and irrelevant correlations under limited data conditions, we performed a targeted experiment. Specifically, we worked with the Heart dataset and randomly augmented it with extraneous columns sourced from the Adult dataset, ensuring these columns held no relevance for the Heart classification task. We assessed model performance as we gradually increased the quantity of these irrelevant columns. To maintain experimental consistency, each method was provided the same set of 8 labeled examples, with no access to additional unlabeled data. Thus, approaches requiring unlabeled data, such as SCARF and STUNT, were excluded from this comparison.

Performance impacts stemming from the introduction of spurious correlations are depicted in Figure~\ref{fig:spurious}. Among all tested methods, \textsf{FeatLLM} exhibited the smallest drop in accuracy due to noisy correlations. This resilience can likely be attributed to our framework’s explicit mechanism for linking features to the target task through prior knowledge when rule extraction occurs. As a result, the framework largely avoids producing rules informed by unrelated or noisy features, enhancing robustness. In practice, we noted that only 1-2\% of extracted rules utilized the introduced noisy columns. Nevertheless, it is important to acknowledge that \textsf{FeatLLM} remains susceptible to learning spurious patterns and misattributing causal significance if columns inserted come from datasets that are strongly associated with the target task (see Appendix~\ref{sec:spurious-appendix}).

\vspace*{-0.12in}
\paragraph{Hyper-parameter analysis}
\textsf{FeatLLM} relies on two central hyper-parameters: the ensemble count and the number of rules retrieved per class within each LLM invocation. To understand how these settings influence outcomes, we conducted a series of experiments. Our findings indicate that while increasing the ensemble number enhances performance, the benefits tend to plateau after reaching approximately 20 ensembles (refer to Fig.~\ref{fig:appendix-hyperparameter1} in Appendix). When it comes to the rule count, our statistical tests did not reveal a significant effect; nonetheless, we selected 10 rules as the optimal compromise between prompt size and predictive performance (see Fig.~\ref{fig:hyperparameter2}). We hypothesize that augmenting the number of rules raises the input dimensionality, which, although it brings more information, also fosters overfitting—especially notable in few-shot scenarios.

\section{Conclusion}
In this paper, we introduce \textsf{FeatLLM}, setting a new benchmark for few-shot tabular data analysis with LLMs. Distinct from approaches where LLMs make direct sample-wise predictions, \textsf{FeatLLM} harnesses LLM capabilities to generate heuristic criteria for prediction, drawing an analogy with feature engineering. Through rule extraction and the use of bagged ensembles, \textsf{FeatLLM} operates with low inference demands and requires no model training—just API access—while also alleviating limitations related to data feature size. Our method consistently demonstrates robust predictive ability and stands out as a data-efficient solution for integrating LLMs into real-world tabular tasks.

At present, \textsf{FeatLLM} is optimized specifically for low-shot learning. Looking ahead, we intend to enhance its ability to design novel features for datasets with more abundant samples, broaden its feature generation beyond simple rules, and emphasize interpretability, thereby enabling its adaptation to a broader spectrum of tasks.

\section{Broader Impact} 
This work introduces \textsf{FeatLLM}, a system that leverages the extensive prior knowledge within LLMs to tackle tabular learning problems and to surpass the standard limitations of supervised tabular learning. By enhancing data efficiency, our research contributes significantly to alleviating the problem of overfitting in scenarios where available labeled data is sparse, using the LLM's pre-existing domain knowledge. \textsf{FeatLLM} offers clear practical advantages for professionals dealing with limited annotated data—common in fields like finance and healthcare, where high labeling costs or acquisition challenges prevail and where the foundational knowledge embedded in LLMs, trained on vast domain-relevant corpora, may be especially beneficial. One crucial line of future inquiry for more extensive application of \textsf{FeatLLM} is to carefully examine the effects of societal biases or potential misinformation present in the LLM’s pretraining data, since such inherited knowledge could substantially influence or even outweigh the information from the limited labeled samples available.

\end{document}